\documentclass[10pt,a4paper]{amsart}
\usepackage[margin=19mm]{geometry}
\usepackage{amsmath,amssymb,mathtools}
\usepackage[scr=rsfs,cal=euler]{mathalfa}
\usepackage{xfrac}
\usepackage[unicode,hidelinks,bookmarks=false]{hyperref}
\newcommand{\ie}{i.e.\ }
\newcommand{\cf}{cf.\ }
\newcommand{\resp}{respectively\ }
\newcommand{\Subsection}{\subsection}
\providecommand{\nobreakdash}{\nobreak\hbox{-}}
\setlength{\emergencystretch}{2em}
\multlinegap0pt
\usepackage{amssymb}
\usepackage{mathtools}
\newtheorem{prop}{Proposition}
\newtheorem{thm}{Theorem}
\newcommand{\E}{\mathbb{E}}
\renewcommand{\P}{\mathbb{P}}
\newcommand{\Z}{\mathbb{Z}}
\title{The 2D Toda lattice hierarchy for multiplicative statistics of Schur measures}
\author{Pierre Lazag}
\date{Source: arXiv:2403.18141v3 (CC BY 4.0)}
\begin{document}
\maketitle

\noindent\emph{Source and licence.} The 2D Toda lattice hierarchy for multiplicative statistics of Schur measures. Version arXiv:2403.18141v3 (CC BY 4.0), distributed by arXiv under the Creative Commons Attribution 4.0 International licence, http://creativecommons.org/licenses/by/4.0/, as recorded in the arXiv licence metadata for this exact version and shown on the abstract page https://arxiv.org/abs/2403.18141v3. Full licence notice: \texttt{licenses/arxiv-2403.18141-CC-BY-4.0.txt}.\par\smallskip

\noindent\textit{Editorial scope.} Counted unit: the article's prop prop:prop2 (source0.tex lines 162--167). The article's complete original proof of it is delivered with it (source0.tex lines 345--386, one \texttt{proof} environment of 42 lines, which the article places away from the statement). No mathematics is written, repaired or re-derived: the statement and the proof are the article's own lines, in the article's own notation, and no retained line is substituted or re-flowed.  Retained uncounted as the source-local context the proof needs: lines 138--140; lines 158--160; lines 266--274; lines 299--301.  Nothing was omitted from any retained span, so the package carries no omission record.  No retained line contains a citation, so the wrapper carries no bibliography and no bibliography entry.  No retained line contains a figure environment, an image-inclusion command or drawing code, so no image is delivered and the package declares no asset dependency.  Editorial formatting only, in addition: an AMS \texttt{amsart} preamble that restates the article's own declarations and macros used by the retained lines, namely the environments \texttt{prop}, \texttt{thm} on separate counters, together with the standard packages those lines need.  The retained environments print in this document as Proposition 1, Theorem 1 in order of appearance, whereas the article prints the same items with the numbering of its own full document.  The complete native source of the article as distributed in the arXiv e-print for this exact version is bundled unchanged as \texttt{sources/156728/source0.tex}.
\begin{align} \label{eq:defK}
K_{t,t',\sigma}(x,y) = \sum_{k \in \Z'} \sigma(k) J_{x +k}(t,t') J_{-y-k}(-t,-t'), \quad x,y \in \Z'.
\end{align}

\begin{align} \label{eq:defKshort}
   K_{t,t'} := K_{t,t', \mathfrak{1}_{\Z'_{>0}}}. 
\end{align}

\begin{thm} \label{thm:okounkov} For any finite set $X= \{x_1, \dots, x_m\} \subset \Z'$, we have
\begin{align*}
\P_{t,t'} ( X \subset \mathfrak{S}_0( \lambda ) ) = \det \left( K_{t,t'} (x_i, x_j) \right)_{i,j=1}^m,   
\end{align*}
where the correlation kernel $K_{t,t'}$  is given by (\ref{eq:defK}) and (\ref{eq:defKshort}):
\begin{align*}
K_{t,t'} (x,y) = \sum_{ k \in \Z', \hspace{0.1cm} k>0} J_{x+k}(t,t') J_{-y-k}(-t,-t').
\end{align*}
\end{thm}

\begin{align} \label{eq:factorKsigma}
    K_{t,t',\sigma} = H_{t,t'}\sigma \check{H}_{t,t'},
\end{align}

\begin{prop} \label{prop:prop2} Let $\P_{t,t'}$ be the Schur measure with parameters $t,t'$.
Then, we have
\begin{align*}
\tau_n(t,t';\sigma) = Z_{t,t'} \E_{\P_{t,t'}} \left[ \prod_{x \in \mathfrak{S}_0(\lambda)} (1- \sigma(x-n) )\right].
\end{align*}
\end{prop}

\begin{proof}[Proof of Proposition \ref{prop:prop2}]
By Theorem \ref{thm:okounkov}, we have
\begin{align*}
    \E_{\P_{t,t'}} \left[ \prod_{x\in \mathfrak{S}_0(\lambda)}(1-\sigma(x-n) )\right] = \det(1- \sigma(\cdot-n)K_{t,t'})_{\ell^2(\Z')},
\end{align*}
so it suffices to prove that
\begin{align*}
    \det(1-K_{t,t',\sigma})_{\ell^2\{n+1/2,n+3/2,\dots\}}= \det(1- \sigma(\cdot-n)K_{t,t'})_{\ell^2(\Z')}.
\end{align*}
Using the factorization (\ref{eq:factorKsigma}), we use of the identity
\begin{align*}
\det(1 + AB ) = \det(1 +BA),
\end{align*}
valid when $A$ and $B$ are Hilbert-Schmidt operators. We obtain
\begin{align*}
\det \left( 1 -K_{t,t'\sigma} \right)_{\ell^2 \{n+1/2,n+3/2, \dots \}} &= \det \left( 1 - \mathfrak{1}_{ \{n+1/2,n+3/2, \dots \} } H_{t,t'} \sigma \check{H}_{t,t'} \mathfrak{1}_{ \{n+1/2,n+3/2,\dots \} } \right)_{\ell^2(\Z')}\\
&= \det \left(1 - \sqrt{\sigma} \check{H}_{t,t'} \mathfrak{1}_{\{n+1/2,n+3/2,\dots\}}H_{t,t'} \sqrt{\sigma} \right)_{\ell^2(\Z')}.
\end{align*}
We can now observe that the kernel of the left factor can be expressed as:
\begin{align*}
    \left(\sqrt{\sigma} \check{H}_{t,t'} \mathfrak{1}_{\{n+1/2,n+3/2\dots\}}\right)(x,y)&=\sqrt{\sigma(x)}J_{x+y}(t',t)\mathfrak{1}_{\{n+1/2,\dots\}}(y) \\
    &= \sqrt{\sigma(x)}J_{x+y}(t',t)\mathfrak{1}_{\Z'_{>0}}(y-n) \\
    &=\sqrt{\sigma(x+n-n)} J_{x+n+y-n}(t',t) \mathfrak{1}_{\Z_{>0}}(y-n)\\
    &=\left(T^n\sqrt{\sigma(\cdot-n)}\check{H}_{t,t'}\mathfrak{1}_{\Z_{>0}}\right)(x,y-n),
\end{align*}
where $T: \ell^2(\Z') \to \ell^2(\Z')$ is the shift operator: $Tf(x)=f(x+1), \, f \in \ell^2(\Z')$,  $x \in \Z'$. Similarly, the kernel of the right factor can be written as
\begin{align*}
    \left( \mathfrak{1}_{\{n+1/2,n+3/2,\dots\}} H_{t,t'}\sqrt{\sigma} \right)(y,x') = \left(\mathfrak{1}_{\Z'_{>0}} H_{t,t'}\sqrt{\sigma(\cdot-n)}T^{-n}\right)(y-n,x').
\end{align*}
It follows that
\begin{align*}
    \left(\sqrt{\sigma} \check{H}_{t,t'} \mathfrak{1}_{\{n+1/2,n+3/2,\dots\}}H_{t,t'} \sqrt{\sigma}\right)(x,x')=\left(T^n\sqrt{\sigma(\cdot -n)}\check{H}_{t,t'}\mathfrak{1}_{\Z'_{>0}}H_{t,t'}\sqrt{\sigma(\cdot -n)} T^{-n}\right)(x,x'),
\end{align*}
and thus
\begin{align*}
\det \left( 1 -K_{t,t'\sigma} \right)_{\ell^2 \{n+1/2,n+3/2, \dots \}}&= \det \left( 1- \sqrt{\sigma (\cdot -n) } \check{H}_{t,t'} \mathfrak{1}_{\Z'_{> 0}} H_{t,t'} \sqrt{\sigma( \cdot - n) } \right)_{\ell^2(\Z')} \\
&= \det \left( 1- \sigma( \cdot - n ) \check{ H}_{t,t'} \mathfrak{1}_{ \Z'_{> 0}} H_{t,t'}  \right)_{\ell^2(\Z')}\\
&= \det\left(1- \sigma(\cdot-n)K_{t',t}\right)_{\ell^2(\Z')}\\
&=\det\left(1- \sigma(\cdot -n)K_{t,t'}\right)_{\ell^2(\Z')},
\end{align*}
where the last equality follows from the symmetry $K_{t',t}(x,y)=K_{t,t'}(y,x)$, which does not affect the value of the determinant. The Proposition is proved.
\end{proof}

\end{document}
